\documentclass{article}
\usepackage{amsmath}
\usepackage{xcolor}
\begin{document}

INTRODUCCION A LAS ECUACIONES DIFERENCIALES



Que es una Ecuacion Diferencial?

Una E.D. es un expresion matematica que tiene presente dentro de si derivadas.

El objetivo, es identificar una función solución que satisfaga dicha expresion


\begin{gather*}
x^{2}+5x+6=0 \\
(-3)^{2}+5(-3)+6=0 \\
x_{1,2} = ??   \rightarrow  x=-3 \\
\text{inicialmente en esta expresion, la idea era} \\
\text{determinar el valor de "}x\text{" como valor real}
\end{gather*}


En una E.D. se va tener presente, derivadas que pueden depender de una o mas variables


\begin{gather*}
\frac{d^{2}y}{dx^{2}}+5\frac{dy}{dx}+6y=0 \\
donde y = variable dependiente \\
x = variable independiente \\
donde el objetivo es encontrar una funcion y=(x) \\
Asumamos que   y = e^{-3x} ,  x^{2}+3  , sen(4x) \\
                                 y =  e^{-3x}  \rightarrow \text{solucion a la E.D.} \\
y' = -3e^{-3x} \\
y'' = 9e^{-3x} \\
9e^{-3x}-15e^{-3x}+6e^{-3x}=0
\end{gather*}


Clasificacion de las E.D.

E.D. Ordinarias $\rightarrow$  tienen presente derivadas simples


\begin{gather*}
\frac{dx}{dt} +5x =0 
\end{gather*}


E.D No Ordinarias $\rightarrow$  tienen presente derivadas parciales


\begin{gather*}
\frac{\partial x}{\partial y} + \frac{\partial x}{\partial z}=3x
\end{gather*}


Segun el orden, grado, linealidad

\textbf{Orden}$\rightarrow$  Se puede identificar por la \colorbox[RGB]{248,231,28}{derivada mayor} presente en la E.D.


\begin{gather*}
\frac{dx}{dt} +5x =0 \rightarrow \text{ Primer Orden} \\
\frac{d^{3}x}{dt^{3}} +\frac{d^{2}x}{dt^{2}}+5x =0  \rightarrow \text{Tercer Orden} \\
\frac{d^{2}y}{dx^{2}}+5\frac{dy}{dx}+6y=0  \rightarrow \text{ Segundo Orden}
\end{gather*}




\textbf{Grado} $\rightarrow$  El exponente mayor presente en la E.D del termino que tiene el orden de la E.D.


\begin{gather*}
(\frac{d^{2}y}{dx^{2}})^{3}+5(\frac{dy}{dx})^{5}+6y=0  \rightarrow \text{Grado 3} \\
       A                  B                      
\end{gather*}




\textbf{Linealidad} $\rightarrow$  Basado en el formato plantilla en forma de escalera *


\begin{gather*}
a_{n}(x)\frac{d^{n}y}{dx^{n}} + a_{n-1}(x)\frac{d^{n-1}y}{dx^{n-1}} + .....a_{1}(x)\frac{dy}{dx}+a_{0}(x)y = g(x) \\
\frac{d^{3}y}{dx^{3}} +e^{3x}\frac{d^{2}y}{dx^{2}}+5x\frac{dy}{dx}-3y =sen(x)  \rightarrow \text{lineal} \\
\vspace{0.5em} \\
\frac{d^{3}y}{dx^{3}} +e^{3x}\frac{d^{2}y}{dx^{2}}+5x\frac{dy}{dx}-3y =\text{\colorbox[RGB]{248,231,28}{$sen$}}\text{\colorbox[RGB]{248,231,28}{$($}}\text{\colorbox[RGB]{248,231,28}{$y$}}\text{\colorbox[RGB]{248,231,28}{$)$}}  \rightarrow \text{no lineal} \\
\frac{d^{3}y}{dx^{3}} +e^{3x}\frac{d^{2}y}{dx^{2}}+5x\text{\colorbox[RGB]{248,231,28}{$($}}\text{\colorbox[RGB]{248,231,28}{$\frac{dy}{dx}$}}\text{\colorbox[RGB]{248,231,28}{$)$}}\text{\colorbox[RGB]{248,231,28}{$^{2}$}}-3\text{\colorbox[RGB]{248,231,28}{$\sqrt{y}$}} =sen(x)  \rightarrow \text{ no es lineal} \\
\text{\colorbox[RGB]{248,231,28}{$y$}}\frac{d^{3}y}{dx^{3}} +e^{3x}\frac{d^{2}y}{dx^{2}}+5x\frac{dy}{dx}-3y =sen(x) \rightarrow \text{ no es lineal} \\
\text{Por ejemplo para simplificar el concepto, podemos arrancar} \\
\text{de una ecuacion de 2do grado de colegio} \\
Ax^{2}+Bx+C=0 \\
5x^{2}-3x+8=0 \\
Ax^{4}+Bx^{3}+Cx^{2}+Dx+E=0
\end{gather*}


E.D. de Variable Separable

E.D. Homogeneas

E.D. Lineales


\end{document}
